% ============================================================
%  CHAPITRE 3 : ANALYSE ET CONCEPTION
% ============================================================
\chapter{Analyse et conception}

\section*{Introduction}
\addcontentsline{toc}{section}{Introduction}
% [Copier depuis Canva]

% -----------------------------------------------------------
\section{Identification des acteurs}
% -----------------------------------------------------------
% [Copier depuis Canva]

% -----------------------------------------------------------
\section{Spécification des besoins}
% -----------------------------------------------------------

\subsection{Besoins fonctionnels}
% [Copier depuis Canva — liste visiteur, client, admin]

\subsection{Besoins non fonctionnels}

\begin{enumerate}
  \item \textbf{Performance} — [Copier]
  \item \textbf{Ergonomie et expérience utilisateur} — [Copier]
  \item \textbf{Compatibilité et Responsive Design} — [Copier]
  \item \textbf{Multilinguisme} — [Copier]
  \item \textbf{Multi-devises} — [Copier]
  \item \textbf{Sécurité} — [Copier — JWT, bcryptjs, helmet, sanitizeBody, CORS, rate limiting]
  \item \textbf{Maintenabilité} — [Copier]
  \item \textbf{Scalabilité} — [Copier]
  \item \textbf{Disponibilité et fiabilité} — [Copier]
\end{enumerate}

\begin{table}[H]
\centering
\caption{Règles de rate limiting par route}
\label{tab:rate_limiting}
\begin{tabularx}{\textwidth}{X c c}
\toprule
\rowcolor{beige}
\textbf{Route protégée} & \textbf{Fenêtre} & \textbf{Requêtes max} \\
\midrule
Global (toutes routes)             & 15 min  & 300 \\
Authentification (login/register)  & 15 min  & 20  \\
Changement de mot de passe         & 15 min  & 10  \\
Formulaire de contact              & 1 heure & 5   \\
Demande de devis configurateur     & 1 heure & 10  \\
Passage de commande                & 1 heure & 20  \\
\bottomrule
\end{tabularx}
\end{table}

\begin{table}[H]
\centering
\caption{Résultat de l'audit npm audit}
\label{tab:npm_audit}
\begin{tabular}{lc}
\toprule
\rowcolor{beige}
\textbf{Niveau} & \textbf{Nombre} \\
\midrule
Critique  & 0 \\
Haute     & 0 \\
Modérée   & 2 \\
Faible    & 0 \\
\bottomrule
\end{tabular}
\end{table}

% -----------------------------------------------------------
\section{Conception fonctionnelle}
% -----------------------------------------------------------

\subsection{Diagramme de cas d'utilisation}

\begin{figure}[H]
  \centering
  \includegraphics[width=\textwidth]{use_case.png}
  \caption{Diagramme de cas d'utilisation}
  \label{fig:use_case}
\end{figure}

\subsection{Description des cas d'utilisation}
% [Tableaux CAS 1 à 4 — copier depuis Canva]

\subsection{Diagramme de séquence}

\begin{figure}[H]
  \centering
  \includegraphics[width=\textwidth]{seq_commande.png}
  \caption{Diagramme de séquence — Passer une commande}
  \label{fig:seq_commande}
\end{figure}

\begin{figure}[H]
  \centering
  \includegraphics[width=\textwidth]{seq_configurateur.png}
  \caption{Diagramme de séquence — Utiliser le configurateur}
  \label{fig:seq_configurateur}
\end{figure}

\subsection{Diagramme d'états-transitions}

\begin{figure}[H]
  \centering
  \includegraphics[width=\textwidth]{etats_transitions.png}
  \caption{Diagramme états-transitions du cycle de vie d'une commande}
  \label{fig:etats_transitions}
\end{figure}

\subsection{Diagramme de classe}

\begin{figure}[H]
  \centering
  \includegraphics[width=\textwidth]{diagramme_classe.png}
  \caption{Diagramme de classe}
  \label{fig:diagramme_classe}
\end{figure}

% -----------------------------------------------------------
\section{Conception générale}
% -----------------------------------------------------------

\subsection{Architecture générale de la solution}

\begin{figure}[H]
  \centering
  \includegraphics[width=\textwidth]{architecture.png}
  \caption{Schéma d'architecture générale}
  \label{fig:architecture}
\end{figure}

\subsection{Processus de conception et de production}
% [Copier depuis Canva]

\subsection{Choix techniques et technologies utilisées}

\begin{table}[H]
\centering
\caption{Outils, technologies et langages utilisés}
\label{tab:technologies}
\begin{tabularx}{\textwidth}{>{\bfseries}p{3.5cm} p{3cm} X}
\toprule
\rowcolor{beige}
\textbf{Élément} & \textbf{Type} & \textbf{Description} \\
\midrule
Visual Studio Code   & Outil           & Éditeur de code utilisé pour le développement frontend et backend. \\
Blender 5.0          & Outil           & Logiciel de modélisation et de rendu 3D pour les visuels photoréalistes. \\
Git                  & Outil           & Système de gestion de versions. \\
GitHub               & Outil           & Plateforme d'hébergement et de gestion du code source. \\
Postman              & Outil           & Outil de test et validation des routes de l'API REST. \\
MongoDB Atlas        & Outil           & Service cloud pour héberger et gérer la base de données. \\
Draw.io              & Outil           & Création des diagrammes UML et schémas d'architecture. \\
\midrule
React 19             & Framework Frontend   & Architecture basée sur des composants réutilisables. \\
Vite                 & Outil de build       & Développement et optimisation du frontend. \\
TailwindCSS          & Framework CSS        & Interface responsive et cohérente. \\
React Router v7      & Framework Frontend   & Gestion de la navigation en mode SPA. \\
i18next + react-i18next & Bibliothèque Frontend & Internationalisation FR/AR/EN avec support RTL. \\
Vitest + React Testing Library & Outil de test & Tests unitaires du frontend. \\
\midrule
Node.js              & Technologie Backend  & Environnement d'exécution JavaScript pour le backend. \\
Express.js           & Framework Backend    & Création de l'API REST. \\
MongoDB + Mongoose   & Base de données      & Gestion des données de l'application. \\
\midrule
JavaScript (ES2022+) & Langage & Backend et écosystème React. \\
TypeScript           & Langage & Sur-ensemble typé de JavaScript pour la qualité du code. \\
HTML5                & Langage & Structure des pages web. \\
CSS3                 & Langage & Mise en forme et animations visuelles. \\
\bottomrule
\end{tabularx}
\end{table}

\subsection{Justification des technologies utilisées}
% [Copier depuis Canva]

\subsection{Schéma de la base de données MongoDB}

\begin{figure}[H]
  \centering
  \includegraphics[width=\textwidth]{schema_mongodb.png}
  \caption{Schéma de base de données MongoDB}
  \label{fig:schema_mongodb}
\end{figure}

\section*{Conclusion}
\addcontentsline{toc}{section}{Conclusion}
Ce chapitre a permis de définir les bases conceptuelles et techniques de la plateforme FENDRI. Les besoins fonctionnels et non fonctionnels ont été identifiés et modélisés à travers les diagrammes UML. L'architecture générale ainsi que les choix technologiques ont été établis et justifiés. Le chapitre suivant sera consacré à la réalisation et à l'implémentation de la solution développée.
